\quellseite{635}
\part{VI. The Close of Religious Instruction}

\quellseite{637}
\chapter[§ 39. THE PRELIMINARY CLOSE]{§ 39. THE PRELIMINARY CLOSE BY CONFIRMATION}

J. W. Hoefling, \emph{Sakr. d. Taufe}, 2 vols., 1859. G. v. Zezschwitz i, pp. 580–726. Th. Harnack, \emph{Katechetik}, pp. 177–196. K. Buchrucker, pp. 96–103. E. Sachsse, pp. 401–418. O. Baumgarten, pp. 81–85. J. Gottschick, pp. 119–124. E. Chr. Achelis, pp. 303–335. J. Steinbeck, pp. 238–249. *W. Caspari, \emph{Die evangelische Konfirmation, vornehmlich i. d. luth. Kirche}, 1890. W. Diehl, \emph{Zur Geschichte d. Konfirmation}, 1897. Ernst and Adam, \emph{Die katechetische Geschichte d. Elsasses}, 1897. E. Simons, \emph{Konfirmation und Konfirmandenunterricht}, 1900. *G. Kawerau, \emph{Bedarf d. gegenwärtige Konfirmationsordnung einer Aenderung?} in \emph{Halte, was du hast}, 1901. E. Chr. Achelis, \emph{Offener Brief an Prof. Kawerau}, ibid., 1901. E. Chr. Achelis, \emph{Die Bestrebungen zur Reform d. Konfirmationspraxis} in \emph{Theol. Rundschau}, 1901 and 1904. G. Beelitz, \emph{Unsere Konfirmationsordnung i. Licht d. Hl. Schrift}, 1901. *E. Hansen, \emph{Die Geschichte d. Konfirmation in Schleswig-Holstein}, 1906. Freie kirchl.-soziale Konferenz, Heft 11. 12. 15. 16. 23. 24. K. Bonhoff, \emph{Die Unhaltbarkeit d. Forderung d. Konfirmationsgeluebdes}, 1908. E. Simons, \emph{Die Konfirmation}, 1908. G. Rietschel, \emph{Lehrbuch d. Liturgik} ii, 1909. *F. Renndorff, \emph{Das Problem d. Konfirmation}, 1910. *E. Sehling, \emph{Evang. Kirchenordnungen d. 16. Jahrhunderts}, 1903 ff. *M. Reu, \emph{Quellen z. Geschichte d. kirchl. Unterrichts}, 1904 ff.

If the Church, by the educational agencies above mentioned, by the use of the material indicated, and in pursuance of a method adapted both to the inner life of the child and to the character of the teaching material, has done everything that pastoral fidelity may suggest, it is likely that she has accomplished everything laid down by us as the aim of her teaching and training, unless the pupil's soul has been actually barred against her efforts (pp. 238 f., 279 f.). True, a life of personal faith may have been aroused only in rare instances; but the sacred truths, upon which the life of the mature congregation is fundamentally based, and by which it is constantly renewed, have been imbedded and anchored in the intellect of the adolescent youth; their emotional life has been stirred to a
\quellseite{638}
vital interest in these truths, and the will has been habituated to the pursuit of those paths in which, soon or late, the Spirit can and will arouse the soul to personal faith. So far as men are a factor in the premises, a general participation in the life of the mature congregation has thus been rendered possible: the member of the Church, having “become of age” (p. 275), is henceforth able to take part in the counsels and activities of the adult congregation, and thus the specific period of teaching and training has come to a close. Before the formal close of this period, however, another important step has been taken, when the congregation, in so far as that is in her power, has passed upon the ripeness of her children for Communion, and thus fully received them into her fellowship of worship. We distinguish accordingly between the preliminary and the final close of religious instruction.

We discuss first confirmation since the preliminary close of religious teaching in the Lutheran Church generally coincides with this act. Through Baptism the children have become members of the Church and, therewith, of the Kingdom of Christ on earth. The prevenient grace of God has, objectively, imparted everything to them which the grace of Christ has acquired. They have become children of God, and He has become their father. This covenant of grace being perpetual, it requires on the part of God neither repetition nor amplification. Notwithstanding we do not admit infants and children to Holy Communion, a refusal which implies rejection on our part of the communio infantium, which was in general vogue in the Early Church and, here and there, also in the Medieval Church until the twelfth century, later to be adopted once more by the Bohemian Brethren (pp. 53, 75, 79; compare Luther's letter of 1523 to Hausmann, and the injunction by the Fourth Lateran Council, of 1215, to grant communion to children who have passed the seventh year). We reject infant communion because Holy Communion is not in
\quellseite{639}
the same sense essential to salvation as the Word and Baptism; and because of the clear direction of Paul in 1 Cor. 11:28. Inasmuch as participation in Communion is conditioned by self-examination, the mature congregation and the divinely instituted office for the administration of the means of grace are enjoined from admitting those to the Holy Supper who lack the ability to examine themselves. The Church, however, has the duty to train her immature members for such self-examination and subsequent participation also in that part of her worship in which, as yet, they have no part. It is accordingly one of the functions of the training given by the Church to the young of her fold, not only to point back again and again to the divine act of Baptism, but also to direct their eyes to the blessed gift upon the altar, and to prepare them for its reception. This is true of catechumenal instruction in particular, which, more than all other teaching, should bear the character of preparation and training for the Sacrament (cp. pp. 22 ff., 89 f., 430); which should make the training of heart and mind its aim, and in which the mature catechist discloses to his youthful friends the importance of each truth for the Christian life more than before. If the required knowledge of sin and salvation, that is to say, the capacity for self-examination, has been attained, the mature congregation has no longer the right to withhold Holy Communion from those hitherto immature. They now join the number of those whose title to Holy Communion and full participation in the worship of the Church must be conceded. With confirmation, the instruction of the young has found its preliminary close. The age at which that is done is in itself irrelevant; but the general (although by no means universal) practise of fixing the transition period from childhood to youth as the time for confirmation, has the strongest arguments in its favor (p. 263). Much can be said in favor of the time from the sixteenth to the seventeenth year (p. 268); in many regions of Germany it has
\quellseite{640}
been given the preference, and some want to see this made a general rule; but who gives the Church the right to withhold Holy Communion longer than necessary from many of her members; and who gives us the assurance that we can reach at this later age the same high percentage of those baptized in infancy?

Should the close of instruction be signalized and authenticated also outwardly by a specific rite to be performed before and by the adult congregation? No absolute necessity for such a rite exists: what we call confirmation is merely an institution of the Church; and the laying on of hands (Acts 8:19; Heb. 6; 1 Tim. 4; 2 Tim. 1, etc.), which has been used as argument is not at all spoken of by Scripture as connected with the close of religious instruction and apart from Baptism. In the very homeland of the Reformation, Electoral Saxony, there existed, down to the eighteenth century, no special solemnity whatever in connection with the close of catechetical instruction. Yet we have this opinion by Luther: “Confirmatio, ut volunt episcopi, non curanda; sed tamen quisque pastor posse scrutari a pueris fidem, quae si bona et germana esset, ut imponeret manus et confirmaret, non improbamus” (Sermon for Laetare, Luther’s Works, Weimar Edition xi, 66). And in the \emph{Reformatio Wittenbergensis} of 1545 (German in Sehling, \emph{Kirchenordnungen des 16. Jahrhunderts} i, 211), composed by Melanchthon and signed by Luther, we read: “Ad hanc consuetudinem (catechismi) sanciciendam prodesset ritus confirmationis, cum, videlicet exacta pueritia, jam firmior aetas seu adolescentia accederet, palam in ecclesia, audienda esset integra doctrinae confessio, et cum interrogatus promitteret constantiam in hac ipsa sententia recitata et in huius ecclesiae suae confessione, manus pastoris ei imponendae essent, et publica precatione petenda mentis et cordis in hoc confitente confirmatio et gubernatio. Haec ceremonia non esset inane spectaculum, ut nunc est episcoporum
\quellseite{641}
ritus, sed profutura esset ad retinendam doctrinae puritatem et propagationem sententiae ecclesiasticae, ad concordiam et disciplinam.” The fear was prevalent that, by the introduction of such a solemn act something of the leaven of the rightly disesteemed Roman confirmation might be domesticated in Lutheran regions. The whole process of instruction is designated as the true biblical confirmation; wherefore, then, a special rite? This is the unmistakable trend of the Order for Churches in Electoral Saxony of 1580. We read: “Let the pastors diligently instruct the people that catechetical instruction is the true Christian confirmation, that is, a confirming of the faith in which the infant had been baptized. Upon this the stress should be laid in these examinations, and the young should be exhorted throughout life to conform to the teaching which they have received. Instead of doing likewise, the Papists substitute a show for confirmation, which by all Christians is avoided and shunned as a mass of superstition and error” (Reu i 2, p. 143). The word “examination” used in this quotation does not by any means refer to some concluding ceremony but merely to the review of the Catechism during Lent, which the young people and the servants were enjoined to attend year after year. And even where, in addition to the general instruction, a specific course was arranged for the benefit of the candidates for the First Communion, a particular closing solemnity was by no means generally deemed necessary. In 1564 we find the following arrangement in force in the Ansbach region: “It is highly commendable that children about twelve years of age, before being accorded the privilege of Holy Communion, should be enrolled as catechumens at a convenient time of the year, for the purpose of receiving, as a substitute for papal confirmation, and in addition to the usual catechization, for several consecutive days or weeks special instruction in doctrine, so that they may have the requisite intelligence (\emph{Verstandshaben}) for a worthy reception of the
\quellseite{642}
Holy Sacrament of the body and blood of our Lord and Saviour. Such examination and exercise takes place here at Onolzbach for the city children on week-days at twelve o’clock between Easter and Pentecost for one hour each day, and for the village children who belong to this parish on Sundays and festivals during the period between Reminiscere and Exaudi at one o’clock, so that they all may receive the Holy Supper together on Pentecost, after each one shall have made his confession privately on the day previous” Reu, i 1, p. 580).

Notwithstanding, it is readily seen why, already in the century of the Reformation, there developed in many places a specific rite for concluding the period of religious instruction. In the last analysis it is not the pastor alone, but the mature congregation to which the decision is left whether the children, “as far as the intelligence is concerned”, as G. Karg of Ansbach fitly expresses it, are worthy to receive the Holy Communion or not. At all events, the congregation has a vital interest in the new accessions that, in any particular year, “as far as their intelligence is concerned”, are declared ready to be added to the number of those already entitled to Communion; and so a public act for the declaration of such readiness seems to be desirable. But it is of the utmost importance that nothing be incorporated in this act which might be looked upon as a factor supplementary to Baptism, or whereby the reception of Holy Communion becomes a matter of constraint, or whereby violence is done to the children’s veracity, or which would tend to bring about a “holy” congregation in the sense of the Donatists,—errors of which abundant traces are found in the history of confirmation.

Inasmuch as the decision whether the knowledge of sin and salvation necessary to self-examination has been acquired is a prerogative of the mature congregation, it is desirable that the children who have been prepared for their First Communion should be presented and publicly examined. This
\quellseite{643}
should be done at a time when a good attendance can be expected, that is, on some Sunday in the forenoon. Timely notice should be given; all should be urged to attend, the church officers, parents, sponsors particularly. The latter should be invited to be present at the opening of the period of instruction as well, in order to have their attention called to its importance, and to solicit their co-operation in making it a time of blessing. This should not be a mere recitation of the chief parts of the Catechism; nor should it lose itself in a mass of unimportant detail, but rather focus itself upon the great fundamental truths confessed by the Evangelical Lutheran Church as they have been combined in the Catechism and, in the form of an unaffected colloquy, furnish the evidence that the children are conversant with the same. It will, of course, be necessary that now some Bible passage, now some hymn stanza, now some section of the Catechism should be called for and recited by the children; but the public examination dare not degenerate to a mere test of the child’s ability and willingness to memorize; the chief fact to be established is that the fundamental truths of Christian doctrine and life have been mastered. While the bashfulness of many children explains the practise of catechist and examiner to go over the material that is to form the subject of the examination, it is clear that a practise which makes the children beforehand acquainted with the questions to be answered defeats the very purpose of the act, and lowers it to the status of a farce.—When the congregation has recognized that the catechumens possess the requisite knowledge of sin and grace, could it do otherwise than remember them in prayer at the throne of grace? They have been remembered during the time of instruction in the general prayer (cf. p. 35); now it is proper that God should be praised for having led them so far; prayer should be made to Him to still further bless them with knowledge and with the true saving faith, as well; that He should therefore strengthen
\quellseite{644}
such faith as they already have and thus make them real and useful members of His body.—The examination and prayer, then, constitute the chief elements of this rite. The examination should be preceded by a discourse—free or read from the liturgy—on the whole act of confirmation or on the meaning, appropriateness, and sanction of the examination. The prayer should be preceded by an exhortation to sincere thanksgiving and earnest supplications. The prayer may be immediately followed by the imposition of hands and the words of blessing accompanying it (\emph{Einsegnung}). The imposition of hands, in the nature of the case, is nothing but an outward symbol of the impartation to the individual of the blessings that had been invoked in the preceding prayer upon all alike. The blessings invoked by the congregation, and the words of blessing spoken by the pastor over the individual, in which he sums up the prayer of the congregation, are by the laying on of hands made a subject of his bodily perception, as it were. In point of form, these words of blessing accompanying the laying on of hands should be a request, or petition, addressed to God.

The order of the act of confirmation, as here outlined, we find in the Pomeranian Liturgy (the confirmation was introduced in Pomerania as early as 1544). The words are: “The usage of Christian confirmation is observed in the Church for the sake of the Catechism and of prayer, so that the beloved youth may be instructed in Christianity, and be tested in regard to their knowledge of the Catechism. They shall not be admitted to the Holy Sacrament without knowing what they do, which would constitute a peril for them and an offense to others; but only after they shall have learned the Catechism, so that prayers may be said over them by the whole congregation in connection with the imposition of hands, and the blessing be pronounced upon them, whereby they will be so confirmed in their Christianity that their Baptism shall be a comfort to them against the very devil, and they be reminded
\quellseite{645}
of the duty to live before God in the true faith, in holiness, and righteousness”. (Sehling, \emph{Kirchenordnungen des 16. Jahrhunderts}, iv, 1911, p. 443). While in the forma confirmationis the examination is not expressly dwelt upon, it is pre-supposed, and a clear reference to it as something already performed is made to it. We read: “And inasmuch as there are some children here who have recited their Catechism, let us present them to the Lord Jesus Christ and pray with all our heart that He keep and confirm them in the truth through the Holy Spirit, so that they become immovable in their Christianity and bear abundant fruit.” Also an \emph{address} or exhortation to the children, which is even liturgically fixed, is already in existence. This form of confirmation is so appropriate as to arouse the desire in us that, if there was to be a rite of confirmation for bringing the period of instruction to a close at all, this particular form might have gained universal vogue.

But is not, in addition to examination and prayer, the \emph{confession of faith} an integral part of the concluding solemnity? Is it not even the most important element? We meet this view in the sixteenth century mainly where the Church was in conflict with the Ana-baptists and the Schwenckfeldians. Occasionally it found expression in such a way that nothing was required save the bare confession of faith. As a case in point, we read in the Order for the Duchy of Liegnitz and Brieg of 1535: “When the children grow in age and grace, they shall be presented by their parents and sponsors to the ministers before the assembled congregation, so that they may make a public confession of their faith, which is to take the place of confirmation” (Sehling, op. cit., iii, p. 436). Then again the confession of faith was followed by a promise willingly to subject oneself to the discipline of the congregation. Butzer succeeded in incorporating this latter feature in the confirmation act; thus in Hesse in 1538 or 1539; in Strassburg not later than 1543, at least in his own church, though only
\quellseite{646}
for a limited period. From these places the practise spread to other regions. Schwenckfeld and the Ana-baptists had reproached the State Churches, which began to be organized since about 1530, that they had no assurance of the subjective holiness of their members since their membership was linked altogether to infant baptism, “all difference between those without and those within thus being eliminated”. This was indeed a great offense to them; for they did not, like Luther, base the holiness of the Church upon the existence and administration in her midst of the Word and Sacraments, but upon the subjective holiness of her members—a view clearly Donatistic, which induced them to see the real task of the Church in disciplining her members into holiness. From those premises but two conclusions were possible: either substitution of adult baptism for that of infants, to be administered only when subjective holiness and membership in the body of Christ have been attained; or amplification of infant baptism by a ceremony in which the children, upon attaining to mental maturity and an actual surrender to Christ, become subject to the discipline of the Church (the latter is expressly recommended by Schwenckfeld as a way out of the difficulty in case the first step does not find favor, Epistolarium ii, p. 126). Either of these ways, it was urged, would lead to a separation of men and thus to the establishment of a “pure” congregation. Butzer, who stood quite as much under the influence of Zwingli, this enthusiast and spiritual kinsman of the Baptists, as under that of Luther, was deeply impressed with Schwenckfeld’s contention. Accordingly, when he was called to Hesse in 1538, in order to check the Baptist movement, which had found there an extremely fertile field, he saw his surest way to success in recommending for adoption a form for confirmation which contained these two features: surrender to Christ, in the form of a confession of faith, on the part of those baptized
\quellseite{647}
in infancy, and subjection to the discipline of the Church. We do not know his views on this matter; it may be that he did not perceive the implication that only such should have part in confirmation as had surrendered to Christ not with the lip alone but also with the heart; or he may have been of the opinion that everyone has surrendered to Christ who has learned his Catechism; or, finally, that he went back to the popish notion that the Church imparts the Holy Spirit to everyone who submits to the imposition of hands; whatever his view, so much is certain that a foreign element was foisted upon confirmation when he incorporated in it the Creed as a surrender to Jesus Christ or to the Triune God, subjection to the discipline of the Church, and the gift of the Holy Spirit through the laying on of hands. It is an element which has grown neither upon the soil of Scripture nor upon that of the Lutheran doctrine concerning the Church and the means of grace, nor upon that of catechetical expediency, but upon that of postulates extremely dubious; it entered, as a foreign body, the Lutheran Church from without.

In the Ziegenhain Order of Church Discipline, completed in 1538 and published in 1539, which is mainly the work of Bucer, we read: “Children who, through catechization, are far enough advanced in Christian knowledge to be permitted to go to the Lord’s Table, shall, at the instance of the elders and ministers, be presented by their parents and sponsors to the pastors, on the occasion of some great festival such as Christmas, Easter, or Pentecost, and at a place appointed for the purpose. There the elders and all other ministers of the Word shall surround the pastor. Then the pastor shall examine the children in regard to the chief articles of the Christian faith and, when they shall have made reply and publicly surrendered themselves to the Lord Jesus Christ, the pastor shall admonish the congregation to ask the Lord of those children for perseverance in their behalf and for an increase of the Holy Spirit, such prayer to be concluded by a collect. Last of all, the pastor shall put his hands upon the children, thus confirming them in the name of the Lord and receiving them into Christian
\quellseite{648}
fellowship. He should also admit them to the Table of the Lord, but not without an admonition to continue in the obedience of faith, and always to receive in good part and to heed faithfully \emph{Christian discipline and reproof from each and every Christian, especially from the pastors}” (Reu I, 21, p. 411).—In the \emph{Church Regulations of Cassel} of 1539, in close agreement with the regulations of Ziegenhain, we find an “Order of Confirmation and the Laying on of Hands”. Here we read in the introduction: “The ministers of the Word shall prepare the children to be received into full membership during the week before that event; and the children themselves shall severally, each one with his own lips, before the ministers of the Word and the elders, confess their faith and surrender to Christ.” The rite proper thereupon begins with the examination of the catechumens, which has been fixed in liturgical form. However, this examination does not purport to be primarily a test, but rather a confession of the state of the children’s hearts and of their surrender to Christ. For we read in the conclusion: “Do you intend to do and observe all this as you have confessed? Answer: Yes, by the help of our Lord Jesus Christ.—When, thereupon, one of the children shall have made his confession in good, clear language, as here prescribed, this question shall be put to the others one by one: \emph{Do you believe and confess; and will you enter the fellowship and obedience of the Church of Christ just as, at this moment, this child has confessed his faith within your hearing and has entered the fellowship of the Church?} It shall be deemed sufficient when the other children shall answer: Yes, by the help of our Lord Jesus Christ. But they should always be earnestly reminded that they stand before the Lord, whom they cannot deceive, so that their yes is yes indeed. And, some time during the week of preparation, each child shall personally make the prescribed confession before the elders and ministers of the Word. However, the test should not consist in mere words smoothly uttered, since even godly and spiritually mature children may not be capable of fluent answers. What should be insisted on is that the faith is comprehended, not that words are smoothly uttered, which is possible for those very ones who understand them least in the true spirit.” What was meant by surrender to the obedience of the Church, may be seen from the answer to the question: “What does such fellowship of the Church imply?” The answer is: “That I practice strict obedience to the Word of God, by hearing it at stated times, when proclaimed by the
\quellseite{649}
appointed servants or ministers of the Church, on Sundays particularly; also, by submitting with due humility to the reproof administered by the elders and any other Christians for sin on my part and making proper amends; moreover, by instructing and restoring my fellow-Christians whom I find on the wrong way, or informing such other good Christians of the matter who in my judgment are competent to help those in error; but, should these, too, prove powerless in the matter, by laying the matter before the common pastors and elders for adjustment; but, if those in error shall refuse to hear the Church in the persons of said pastors and elders, and are put in the ban in consequence, by likewise treating them as excommunicated, or heathen people”. Such surrender to Christ and passing into the fellowship of the Church is followed by the prayer of the congregation (that prescribed in the Cassel liturgy later passed into many other liturgies, for instance, Loehe’s, although usually somewhat altered). Then follows the laying on of hands: “Thereupon the pastor shall lay his hands upon them and say: ‘Receive the Holy Spirit, as a protection and shield against all evil, as a power and help for all good, from the gracious hand of the Father and the Son and of Holy Spirit. Amen’”. This order of confirmation passed later, in all its fundamental features, into the Church Regulations of 1566 and also those of 1574, which were in force all over the landgraviate of Hesse (Reu, I, 21, pp. 1078—1083). In the enlarged edition of his Catechism of 1537 (issued in 1543, if not earlier) “for the pupils and other children of Strassburg”, Butzer inserted a special section on “Confirmation in the Christian Faith”. Here, after the “Explanation of the Divine Services” and the “Explanation of the Sacraments”, we come upon the following series of questions: “Teacher: What other usages and rites are found in the Christian congregation? Answer: In the first place, the public confirmation of the Christian faith of those baptized in infancy. Teacher: How is that to take place? Ans.: When we children have been instructed in the Christian faith, we are to make confession of it before the whole Christian congregation with our own heart and lips, and thus enter into the covenant with the Lord, and the fellowship and obedience of the Church; having been unable to do this at Holy Baptism, we are confirmed in these things by the Church. Teacher: Whence have you learned this? Ans.: The Scriptures teach that all believers shall enter into the covenant of the Lord by their own faith and confession; and Saint Paul says, Rom. 10:
\quellseite{650}
He that believeth with the heart, etc. Teacher: How shall you be confirmed by the Church in the faith and covenant of the Lord? Ans.: Through the joint blessing and intercession of the Church and the laying on of hands. Teacher: Where do you find authority for these things? Ans.: In the invitation of the Lord Jesus: Suffer the little children to come unto me, etc. Teacher: What does the sign of the laying on of hands signify? Ans.: That the children therewith have been placed under the almighty hand of Almighty God, whose care, protection, and gracious guidance have thus been promised and pledged to them. Teacher: Who is to put his hands upon the children and bless them? Ans.: The regular minister of the Church in behalf of the whole congregation, which is to be a witness to this. Teacher: But the Lord has not enjoined the use of this sign? Ans.: Inasmuch as He and the holy apostles have used this sign for such purposes with great blessing, it behoves us likewise to use it in His name; the Lord will be present with His Spirit and work, and, according to the prayer of the Church, uttered in obedience to the Word, He will graciously confirm such children as members of His kingdom. Teacher.: What is all this to lead to, so far as you are concerned? Ans.: To this that I diligently attend the children's class, and properly learn the Christian faith, confess it in due time and be confirmed therein to my salvation. Teacher: What else? Ans.: To this that I may be assured of the protection and guidance of the divine hand for the purpose of a Christian life in every respect, and for my protection against the old wicked foe, and that I eternally praise and bless God my heavenly Father, through our Lord Jesus Christ. Teacher: In the third place? Ans.: To this that I make every effort to keep the covenant of the Lord, remain in constant touch with His Church, and faithfully receive and use the divine doctrine, the Holy Sacraments, and church discipline”. Butzer’s meaning of penitential discipline (\emph{Busszucht}) is then explained in a special voluminous section in the sense of “obedience of the Church” spoken of by the Cassel Church Order of 1539; an explanation of “covenant with God and true fellowship of the Church of Christ” is afterwards also given in the same Catechism (Reu I, 1, p. 7. 93 ff, 104 f.).

It is to be regretted that, in deference to this precedent, also in other districts those two parts—surrender to Jesus Christ by the confession of faith and submission to the discipline of the Church—were incorporated in the confirmation
\quellseite{651}
rite (cp. especially Sarcerius, who, during his sojourn in Nassau, had probably become conversant with the Church Regulations of Cassel: Reu, I, 21, p. 100 ff.). If the latter element was disregarded, at least the confession of faith was appropriated. This resulted in the assembling of the following three constituent parts of the concluding rite of religious instruction: exploratio, or examinatio; confessio; oratio. These three elements are met with again, for instance, in the Church Regulations for Brunswick-Wolfenbuettel of 1569, where, however, they are less weighted down by Butzer’s notions than is generally the case elsewhere.

Here we read: “When a number of children are found who have become thoroughly acquainted with the chief parts of the Christian faith through their Catechism, a list of them shall be prepared by the pastor of each place and laid before the superintendent at his appearance for the annual visitation, in order that they may be publicly confirmed and strengthened in their faith on some particular Sunday, or on whatever occasion may prove expedient. While the superintendent is present, the bells shall be rung according to custom and the people be summoned together. But the children that have been presented to the superintendent, together with their parents and sponsors, shall stand apart by themselves, preferably before the front altar, where they may be seen by everyone.” After a discourse by the superintendent upon the special purpose of the service “the superintendent shall \emph{examine the children on the chief parts of the Christian faith, requiring them to recite the whole Catechism without the explanations}, whereupon \emph{he shall question them on the chief articles as expressed in Dr. Luther’s Catechism}. When he has found that the answers and confession of faith made are creditable to children of their age, he shall commend their diligence before the whole congregation; he shall dwell at length upon the great treasure possessed by the children in such knowledge, that God has joined Himself to them as their father, and would never, whatever their distress, permit them to perish, as long as they shall persevere in such knowledge, faith, and confession. Then the superintendent shall ask the children whether they will persevere in such knowledge, faith, and confession, and whether, having renounced the devil and all his works and all his ways, they are determined henceforth to abide in this godly endeavor to the end. And they shall answer: Yes, by the grace of almighty God, which we
\quellseite{652}
desire and pray for with all our hearts. Then the superintendent shall once more address words of instruction and exhortation to the congregation assembled, after having given thanks to God for giving to these children His Holy Spirit, by whose grace they have come to a saving knowledge of His dear Son, a fact which should prompt the whole congregation to walk before those children blameless and without offense, and be preserved in such godly conversation unto eternal life”. The superintendent’s exhortation is followed by a prayer of the congregation in behalf of the children (“and wherever the laying on of hands is in favor as a free adiaphoron, there shall be no impediment thereto”), whereupon the superintendent shall conclude the whole ceremony with the blessing of Aaron, and the congregation shall sing “Come, Holy Ghost”, etc. This shall be followed by Holy Communion. It would not only be highly proper, but a further encouragement of the children in their godly endeavor if the parents and sponsors should join them at once in the reception of the Holy Communion” (Richter, \emph{Kirchenordnungen} ii, p. 320 f.).

By no means on a level with these orders of confirmation is Luther’s explanation of 1523 (p. 618); for, while there a confession of the faith of the heart is dealt with, Luther speaks of the fides quæ creditur, as we see from the added clause: “quæ si bona et germana esset”. He confesses his readiness to favor the adoption of a rite which brings out that those baptized in infancy have obtained the knowledge of the true way of salvation, and in which the hands are laid upon those who give evidence of this. What Luther calls scrutatio is simply the examination. A similar examination is also suggested in the \emph{Reformatio Wittembergensis} of 1545 (cf. p. 618), as may be seen from its definition of confirmation as integra doctrinæ confessio, and by its statement of aim: “Profutura esset ad retinendam doctrinæ puritatem et propagationem sententiæ ecclesiasticæ, ad concordiam et disciplinam”. The promise, also, which is to be exacted according to Luther’s recommendation, embraces nothing but a statement of purpose to adhere to the doctrine here confessed. A similar promise appears to have been in the mind of those responsible for the Church Regulations for Pomerania of 1569
\quellseite{653}
(p. 622), where at the end of the exhortation (there is no hint of a question) the assembled children are told: “To this, all say, yes” (Sehling, iv, p. 444). Such a close, indeed, would not be out of keeping with our postulates. As we have no evidence that God, during the time of instruction, intends to work saving faith in all the children entrusted to us and as experience proves, that in many cases he actually does not work it during that period, we absolutely refuse to set the aim of catechetical instruction in the establishment of saving faith (p. 279); how could we now expect at the time of confirmation from all catechized children alike the confession of their saving faith, the faith of their hearts, and require all to go to Holy Communion? Should we not, in that case, do violence to our own better judgment and put a falsehood upon the lips of those children in whom saving faith has not been aroused as yet? No doubt, congregational conditions in a free Church are much better than those in a State Church, and many reasons which require in the latter the abolition of a confession of an inner faith either do not exist among us at all or they do not have the weight which they have elsewhere; but this much is true of the free Church as well that we have no absolute assurance that all the children instructed by us are in a state of saving faith at the end of the period of instruction. For this reason the Church—for the individual cannot simply set aside church order—should no longer hesitate to place the question as to the utility of the present form of confirmation upon the order of the day, and to eliminate from her liturgic forms the confession of an inner faith, especially in view of the fact that she is neither able nor willing to investigate the inner condition of her children and to reject from confirmation all those in whose hearts, in her opinion, saving faith has not yet been enkindled, and thus to create, among the number of the baptized, a communion of souls truly sanctified. Everything that cannot be expected
\quellseite{654}
of all participants should be eliminated from the rite. Quite another thing than the confession of an inner faith is the assent to the correctness of the evangelical way of salvation, especially as formulated in Luther’s Catechism. Such assent, in normal conditions, particularly in a free Church, can be looked upon as a fruit of religious instruction within the reach of all children. If it were solemnly proclaimed in the rite of confirmation, it would virtually be the solemn climax of the examination. With that the promise could quite properly be connected to remain true to the way of salvation recognized as right and to the Church that points it out.\footnote{This dare not be confounded with the pledge peculiar to the confirmation of the Pietists, which arises from different premises, and has a different meaning. Infringing upon Baptism (confirmation=completion of regeneration; renewal of the baptismal covenant on the part of God, etc.) it shows rather an inner affinity to the Hessian Church Regulations. It should be remembered that it was a village in Hesse in which Spener, for the first time, witnessed confirmation (p. 134 f.).} We could, however, refrain from supplying confirmation with this solemn double climax and confine ourselves simply to examination and prayer. But the complete abolition of such a concluding solemnity and, therewith, a return to early usage, as in Electoral Saxony especially (p. 619), would hardly be advisable at this late day. For many a code of Church Regulations has said in truth that it is much easier to convince both children and parents of the necessity of faithfully submitting to instruction when young, if it be made to issue in a solemn rite connected with the solemn presentation and examination of the catechumens before the assembled congregation—a rite in which also the public declaration of the child’s intellectual readiness for the Holy Supper and admission to full participation in the worship of the mature congregation, inclusive of the privilege of sponsorship, may be made,—than if such concluding solemnity should be dispensed with.

When the \emph{privilege of participation in Holy Communion} is given in connection with the rite of confirmation, it can, of
\quellseite{655}
course, have no other force than that of the declaration that the children are now possessed of the intellectual maturity for self-examination (“intellectually prepared” worthily to receive the Holy Supper, George Karg, 1565, p. 620). In no wise should it be construed as affirming the inward worthiness of the catechumens. Much as the catechist should endeavor to impress upon his catechumens throughout the period of instruction, during the last weeks particularly, the blessing of the Holy Supper, and, so far as that is in his power, to arouse a desire for it, he will now, instead of simply declaring his catechumens prepared for the Holy Supper, emphasize very strenuously that this is the case only when repentance for sin, faith in the grace of God, desire for the assurance of forgiveness, dwell in their hearts. He will, furthermore, carefully refrain from anything that might be felt as constraint. For this reason the first communion will not be set for the day of confirmation but for a day a week or two afterward, preferably for a day on which the congregation has been wont to celebrate the sacrament and on which rather a large than a small number of communicants is likely to appear at the altar. He will assign to the catechumens no particular place which would make it easy for the congregation to determine the number of catechumens participating in the sacrament, or even to challenge such investigation. He will rather, whenever the occasion calls for it, make it plain that confirmation does not under all circumstances require participation in Communion. As long as the order of worship observed in the congregation provides for a confession of an inner faith, he will endeavor to impart to children as well as parents the information that no confession of saving faith is forced from anyone on the ground of his submission to the examination. Much as he will endeavor in his capacity as a faithful pastor to guard his catechumens against a false profession and participation in the Sacrament by constraint, he will subsequently be impelled by
\quellseite{656}
the same pastoral fidelity to continue his solicitude for those who have stayed away from the Table of the Lord. He knows who they are; for, if at any time, it is in connection with the first participation in the Holy Sacrament that personal announcement should be insisted upon, and therefore pastoral fidelity will impel him to keep an eye upon them in connection with the various educational channels that continue to flow after confirmation. He trusts in God for the arrival of “time and hour” also for these; but he will beware of dragging it in artificially. If he can only succeed in maintaining the connection of such as these with the Word as means of grace, he awaits further developments in patience and confidence. And should even this connection be torn apart, he still knows that the power of God is not diminished, but that, even in that case, He can lead them back again to the fold. Nor will he forget that such a passage as 2 Cor. 2: 15. 16 is found in the Bible.

It is the duty of the pastor performing the confirmation rite to give to those confirmed a certificate of their confirmation, which, should they change their domicile, would serve for identification. It is advisable that this certificate should be adorned with a Bible passage throbbing with life and power, light and warning—a message from on high whenever and wherever beheld.

\quellseite{657}
\chapter[§ 40. THE FINAL CLOSE]{§ 40. THE FINAL CLOSE OF RELIGIOUS INSTRUCTION}

E. Sachsse, p. 423. *K. Knoke, \emph{Recht und Pflicht der evangelischen Kirche hinsichtlich der relig. Unterweisung ihrer heranwachsenden Jugend}, 1912.

That religious instruction cannot have come to a final close with confirmation, performed at an age of 13 or 14 years, has already been shown in the preceding chapters. We have seen that religious instruction must continue because its aim is not (intellectual) maturity for the participation in the Sacrament, but maturity for full participation in all privileges and duties of the congregation. The educational agencies at our command for such further religious education, and the material available for the purpose have likewise been dealt with. The only question remaining is whether the instruction given after confirmation is to be concluded by a special ceremony. It has been proposed that when the pupil passes from the Junior, or obligatory, branch of the Luther League to the Senior branch, which is optional, that is, at the age of sixteen or seventeen, dismissal from these institutions should be signalized in the form of a special service, accompanied by a certificate showing completion of the post-confirmation course in religious instruction—something like a confirmation certificate—and by conferring upon him at this time, not at confirmation, the privilege of sponsorship. The history of confirmation proves that such a concluding ceremony would tend to make a post-confirmation course in religious training and discipline more general. That children of thirteen are not usually competent to fulfil the obligations assumed in connection with sponsorship, is self-evident; the question whether the privilege of sponsorship might not well be detached from confirmation,
\quellseite{658}
should be seriously considered. So far as the Senior League is concerned, there is less need for a concluding solemnity after several years of membership since identification with it is optional in any event, and a definite time for membership cannot be fixed. However, in the case of men, such a concluding solemnity might coincide with the attainment of their majority and simultaneous reception to voting membership; in the case of women, generally with entrance upon matrimonial life. No other features would be necessary than the gift of a Bible, a prayer, and the benediction. More important, of course, than any concluding solemnity is the further care for all the confirmed. A close of religious instruction and training in the sense of perfect spiritual maturity does not exist here below. While particular forms of the care for souls have come to an end, the general care through the Word and Holy Sacrament is to continue.
